\chapter{Methodology}
\label{ch:method}

\section{Problem statement}
\label{sec:problem}

Let $\src$ be the source domain (SisFall) and $\tgt$ the target domain (Daphnet). Each sample is a
window $x\in\mathbb{R}^{T\times C}$ of tri-axial acceleration and a binary label $y$. In $\src$,
$y=1$ marks a window centred on a fall impact; in $\tgt$, $y=1$ marks a window in which the patient
is freezing. Both labels describe a \emph{motor event that a monitoring system should flag}, which
is the premise of the transfer question. We train $f_\theta$ on labelled samples from $\src$ only and
evaluate it on $\tgt$ without updating $\theta$ (zero-shot transfer). Unlabelled target windows may
be used to estimate input normalisation statistics, since such data are available in any deployment.

The question is not only whether $f_\theta$ transfers, but \emph{why} it does or does not. We
therefore separate three effects:
\begin{itemize}
  \item \textbf{covariate shift}: the distribution of normal movement differs ($p_\src(x\mid y{=}0)\neq p_\tgt(x\mid y{=}0)$), because of age, pathology and sensor placement;
  \item \textbf{label shift}: the prevalence of positives differs (2.6\% of SisFall test windows, 10.0\% of Daphnet windows);
  \item \textbf{concept shift}: the positive class itself is a different phenomenon ($p_\src(x\mid y{=}1)\neq p_\tgt(x\mid y{=}1)$, and the direction that separates the classes in $\src$ need not separate them in $\tgt$).
\end{itemize}
Threshold-free metrics (AUROC) are insensitive to label shift, so a drop in AUROC points to covariate or
concept shift. Normalising the target inputs with target statistics removes first- and
second-order covariate shift per channel; whatever remains after that is attributed to the
representation or to the concept.

\section{A common signal representation}
\label{sec:representation}

Every confounder that is not part of the scientific question is removed by construction
(\Cref{fig:pipeline}).

\paragraph{One sensor, one site.} SisFall carries two accelerometers and a gyroscope on the same
waist board; Daphnet has accelerometers at the ankle, thigh and lower back. We keep the SisFall
ADXL345 and the Daphnet trunk sensor. The two are both worn at the lower back/waist, so the
anatomical site is matched. The second SisFall accelerometer is redundant, and the gyroscope has no
counterpart in Daphnet.

\paragraph{Units and rate.} SisFall counts are converted to $g$ with the datasheet factor
$2\cdot 16/2^{13}$; Daphnet values (mg) are divided by 1000. SisFall is resampled from 200 to
\SI{64}{\hertz} with a polyphase filter that includes anti-aliasing. We chose to downsample the source
rather than upsample the target: interpolated samples carry no information and would have tripled
the cost of every recurrent model.

\paragraph{Anatomical frame.} Each accelerometer is rotated into a vertical (V), antero-posterior
(AP), medio-lateral (ML) frame, and the magnitude $|a|$ is appended, giving $C=4$ channels. For
SisFall the axes are identified from the data instead of a figure in the dataset paper: the
vertical axis carries $\approx 1\,g$ of static acceleration during walking, and after a fall the
subject lies still, so gravity ends up on the AP axis for forward and backward falls and on the ML
axis for lateral ones. The sign convention makes each axis read $+1\,g$ when it points upwards, so a
subject lying prone reads $a_{AP}\approx-1\,g$. Daphnet's trunk axes are documented as forward,
vertical and lateral.

\paragraph{Windows.} Windows last \SI{2}{\second} ($T=128$). Evaluation windows are cut on a \SI{1}{\second}
grid (50\% overlap) in both domains. In Daphnet, samples annotated 0 (outside the protocol) and
sensor dropouts (trunk reading exactly $(0,0,0)$\,mg) are removed, and windows are cut inside each
contiguous valid segment only, so that no window joins two non-adjacent pieces of recording.

\section{Event-level labels}
\label{sec:labels}

SisFall provides one label per file. Labelling every window of a fall file as ``fall'' marks the
walking that precedes the fall and the lying that follows it as positives, and inflates the positive
class to roughly a third of all windows. We locate the impact as the maximum of $|a|$ in each
fall trial (ignoring the first and last second) and define the event interval $E=[t_p-1\,\mathrm{s},\,
t_p+1\,\mathrm{s}]$. A window $w$ is
\begin{equation}
y(w)=\begin{cases}
1 & \text{if } |w\cap E|\geq 1.5\,\mathrm{s},\\
0 & \text{if } |w\cap E| = 0,\\
\text{discarded} & \text{otherwise,}
\end{cases}
\end{equation}
so a positive window always contains the impact with at least half a second of context on each
side, and partially overlapping windows are left out rather than forced into either class. In the
training set only, fall windows are also cut on a \SI{0.25}{\second} grid, which shows the network each
impact at four temporal offsets; validation and test sets keep the \SI{1}{\second} grid. A Daphnet window
is positive when at least half of its samples are annotated as freezing.

\section{Robust, domain-independent scaling}
\label{sec:scaling}

Each channel is standardised as $(x-\mathrm{median})/\mathrm{IQR}$. Fall impacts reach several $g$
while most data sit near $1\,g$, so mean and standard deviation would be dominated by a small number
of windows; median and inter-quartile range are not. Source statistics are fitted on the SisFall
training subjects only. For the target we compare three options, all label-free:
\begin{enumerate}
  \item \emph{source statistics}: Daphnet scaled with the SisFall training statistics (naive deployment);
  \item \emph{domain statistics}: statistics fitted on all Daphnet windows, independently of SisFall (default);
  \item \emph{subject statistics}: statistics fitted on each patient separately, a simple form of test-time adaptation.
\end{enumerate}

\section{Architectures}
\label{sec:arch}

All four models map a window $x\in\mathbb{R}^{128\times 4}$ to one logit and expose an embedding
(the input of the final linear layer) used in the shift analysis. Their sizes are tuned with the
same budget (\Cref{sec:hpo}); the selected configurations are listed in \Cref{tab:hparams}.

\paragraph{CNN-1D.} Three convolutional blocks (convolution, batch normalisation, ReLU), with channel
widths $w,2w,4w$ and kernels narrowing with depth, max-pooling between blocks, global average pooling
and dropout. Global pooling makes the detector indifferent to where in the window the event occurs.

\paragraph{LSTM and GRU.} One or two stacked recurrent layers, optionally bidirectional. The
embedding concatenates the last hidden state with the temporal mean of all hidden states, which
protects against an event early in the window being partly forgotten by the time the sequence ends.

\paragraph{Patch Transformer.} The window is split into non-overlapping patches of $P$ samples
($P\in\{8,16\}$); each patch of all four channels is linearly projected to $d$ dimensions, a learnable
\texttt{[CLS]} token and learned positional embeddings are added, and the sequence goes through $L$
pre-norm encoder layers with four attention heads and GELU feed-forward blocks
\parencite{vaswani2017attention,nie2023patchtst}. The \texttt{[CLS]} output feeds the classifier.

\section{Training protocol and fairness}
\label{sec:training}

For a given seed, every architecture sees the same training windows, in the same mini-batch order,
with the same augmentation draws, because shuffling and augmentation use one seeded generator.
\begin{itemize}
  \item Loss: binary cross-entropy with positive-class weight $n_{\mathrm{neg}}/n_{\mathrm{pos}}$.
  \item Optimiser: AdamW \parencite{loshchilov2019adamw}, weight decay $10^{-4}$, batch size 256, one warm-up epoch followed by cosine decay, gradient clipping at 1.0.
  \item At most 15 epochs; the checkpoint with the lowest (unweighted) SisFall validation cross-entropy is kept, with a patience of 4 epochs. Validation AUPRC saturates at 1.0 after two or three epochs and can no longer rank checkpoints (\Cref{sec:selection}).
  \item Augmentation, in physical units before scaling: rotation about the vertical axis ($\pm20^\circ$), random ML sign flip, amplitude scaling (0.9--1.1) and Gaussian jitter (\SI{0.01}{g}) \parencite{um2017augmentation}. These mimic belt-mounting differences, not pathology.
  \item Decision threshold: the value that maximises $F_2$ on SisFall validation, reused unchanged on the SisFall test subjects and on Daphnet. $F_2$ weighs recall four times more than precision, reflecting the clinical cost of a missed event.
  \item Seeds 0, 1 and 2; results are reported as mean $\pm$ standard deviation.
\end{itemize}

\section{Hyper-parameter search}
\label{sec:hpo}

Each architecture receives six configurations sampled at random from its own search space
(\Cref{tab:space}), trained for at most five epochs (patience 2) on the same training subset (all fall
windows and a fixed half of the non-fall windows). Configurations are ranked by SisFall validation
AUPRC rounded to three decimals, with ties broken by the validation cross-entropy. The Daphnet AUROC of
every configuration is logged but \emph{not} used for selection; it only serves to report the
gain an oracle with access to target labels would obtain \parencite{gulrajani2021domainbed}.

\begin{table}[ht]
\centering
\small
\caption{Hyper-parameter search spaces (six random draws per architecture).}
\label{tab:space}
\begin{adjustbox}{max width=\textwidth}
\begin{tabular}{ll}
\toprule
Architecture & Search space \\
\midrule
CNN-1D & width $w\in\{16,32,48\}$, kernel $\in\{5,7,9\}$, dropout $\in\{0.1,0.3,0.5\}$ \\
LSTM / GRU & hidden $\in\{32,64,96\}$, layers $\in\{1,2\}$, dropout $\in\{0.1,0.3\}$, bidirectional $\in\{$no, yes$\}$ \\
Transformer & $d\in\{32,64,96\}$, layers $\in\{2,3,4\}$, patch $\in\{8,16\}$, dropout $\in\{0.1,0.2\}$ \\
All & learning rate $\in\{5\cdot10^{-4},10^{-3},3\cdot10^{-3}\}$ \\
\bottomrule
\end{tabular}
\end{adjustbox}
\end{table}

\section{Evaluation}
\label{sec:evaluation}

\paragraph{Metrics.} AUROC and AUPRC (threshold-free; AUPRC is more informative under class imbalance
\parencite{saito2015prc}); at the validation threshold: recall (sensitivity), specificity, false-positive
rate, precision, $F_1$, $F_2$ and balanced accuracy; and recall at 10\% false-positive rate. For the
seed-averaged scores, 95\% confidence intervals on AUROC come from a cluster bootstrap that resamples
\emph{subjects}, since windows from one person are not independent.

\paragraph{Experiments.}
\begin{enumerate}
  \item[E1] In-domain fall detection on eight unseen SisFall subjects.
  \item[E2] Zero-shot transfer to all Daphnet recordings, under the three target normalisations.
  \item[E3] Target-domain references with 5-fold patient-grouped cross-validation on Daphnet: training
  from scratch and fine-tuning from the SisFall weights (learning rate $\times0.3$), with six or two
  training patients (the two with the most FoG windows), same training protocol as above. For fold
  $k$, the validation patients (checkpoint selection, threshold) are those of fold $k{+}1$. The
  patients are paired so that every test fold contains FoG: \{S01, S04\}, \{S02, S10\}, \{S03, S05\},
  \{S06, S07\}, \{S08, S09\}.
  \item[E4] Shift analysis on the penultimate-layer embeddings of the SisFall models: squared MMD with an
  RBF kernel \parencite{gretton2012mmd}, the AUROC of a linear domain probe, t-SNE
  \parencite{vandermaaten2008tsne}, and the cosine between the class-difference vectors
  $\Delta_\src=\mu_{\mathrm{fall}}-\mu_{\mathrm{ADL}}$ and $\Delta_\tgt=\mu_{\mathrm{FoG}}-\mu_{\mathrm{walk}}$.
  A negative or near-zero cosine means the direction that the network uses to flag falls does not
  point towards freezing: a direct, representation-level signature of concept shift.
\end{enumerate}
